% ============================================================
%  Subfile: Skema Ontologi (figur + tabel kelas/relasi/properti)
%  Included by bab-3.tex via \subfile{onthology}
% ============================================================
\documentclass[../../thesis]{subfiles}
\begin{document}
\begin{figure}[htbp]
  \centering
  \adjustbox{width=\linewidth, totalheight=0.92\textheight, keepaspectratio, center}{%
  \begin{tikzpicture}[
  node distance=2cm and 2.5cm,
  % Node Styles
  class/.style={circle, draw=blue!80!black, fill=blue!10, thick, 
    minimum size=2.6cm, align=center, font=\bfseries\scriptsize},
  target/.style={circle, draw=blue!80!black, fill=blue!5, dashed, thick, 
    minimum size=2.6cm, align=center, font=\bfseries\scriptsize},
  prop/.style={rectangle, draw=green!60!black, fill=green!10, thick, 
    minimum height=0.5cm, font=\tiny},
  dtype/.style={rectangle, draw=orange!80!black, fill=orange!10, thick, 
    minimum height=0.5cm, font=\tiny},
  % Edge Styles
  edge/.style={->, >={Stealth[length=2.5mm]}, thick, draw=black!70, font=\scriptsize},
  rel/.style={->, >={Stealth[length=2.5mm]}, thick, draw=blue!70!black},
  cross/.style={->, >={Stealth[length=2.5mm]}, thick, dashed, draw=purple!80!black}
]

  % --- 1. HIERARKI KELAS UTAMA (Backbone vertikal, kompak) ---
  % Grade/Subject sebagai akar di atas; Concept (+ ConceptTarget) di bawah.
  \node[class] (grade) {Grade/\\Subject};
  \node[class, below=1.2cm of grade]     (chapter)  {Chapter};
  \node[class, below=1.2cm of chapter]   (subtopic) {Subtopic};
  \node[class, below=1.2cm of subtopic]  (concept)  {Concept};

  \node[target, below=2cm of concept] (target) {ConceptTarget\\(Placeholder)};

  % Relasi Hierarki (arah akar->daun, label di kanan tulang punggung)
  \draw[edge] (grade)    -- node[right] {HAS\_CHAPTER}  (chapter);
  \draw[edge] (chapter)  -- node[right] {HAS\_SUBTOPIC} (subtopic);
  \draw[edge] (subtopic) -- node[right] {HAS\_CONCEPT}  (concept);

  % --- 2. RELASI PADA CHAPTER (loop di kiri) ---
  \draw[edge, draw=orange!90!black] (chapter) edge[loop left, looseness=4.5]
    node[left, align=center, fill=white, inner sep=1pt, font=\tiny]
      {NEXT\_CHAPTER \\ PRASYARAT \\ MEMPERSIAPKAN} (chapter);

  % Summary di kanan Chapter
  \node[prop, right=1.2cm of chapter] (summ) {summary};
  \node[dtype, right=0.3cm of summ] (d_str2) {string};
  \draw[edge] (chapter) -- (summ); \draw[edge] (summ) -- (d_str2);


  % --- 3. DATA PROPERTIES PADA CONCEPT (di kanan) ---
  \node[prop, right=1.5cm of concept] (glos) {glossary\_validated};
  \node[dtype, right=0.3cm of glos] (d_bool) {boolean};
  \draw[edge] (concept) -- (glos); \draw[edge] (glos) -- (d_bool);

  \node[prop, above right=0.6cm and 1cm of concept] (desc) {description};
  \node[dtype, right=0.3cm of desc] (d_str1) {string};
  \draw[edge] (concept) -- (desc); \draw[edge] (desc) -- (d_str1);


  % --- 4. RELASI SEMANTIK (WITHIN-BOOK) — Concept ke Target (di bawah) ---
  \draw[rel] (concept) to[out=-70,  in=70] (target);
  \draw[rel] (concept) to[out=-90,  in=90] (target);
  \draw[rel] (concept) to[out=-110, in=110] (target);

  % Label di kiri celah Concept--Target
  \node[font=\tiny, align=center, fill=white, inner sep=3pt, draw=blue!30,
        rounded corners, anchor=east]
    at ([xshift=-0.6cm]$ (concept)!0.5!(target) $) {
    \textbf{15 Relasi Semantik (\textit{Within-Book}):} \\
    MENDEFINISIKAN, MENYEBABKAN, MEMUNGKINKAN, \\
    MENGATUR, BAGIAN\_DARI, TERDIRI\_DARI, \\
    BERGANTUNG\_PADA, BERINTERAKSI\_DENGAN, \dots \\
    \textit{(Daftar lengkap merujuk pada Tabel)}
  };


  % --- 5. RELASI LINTAS-BUKU (CROSS-BOOK) — loop + label di kiri Concept ---
  \draw[cross] (concept) edge[loop left, looseness=5] (concept);
  \node[left=1.4cm of concept, fill=white, inner sep=2pt,
        font=\tiny, align=left, draw=purple!30, rounded corners] (xblabel) {
      \textbf{Relasi Lintas-Buku:} \\
      LINTAS\_BUKU\_SAMA\_DENGAN \\
      LINTAS\_BUKU\_APLIKASI\_DARI \\
      LINTAS\_BUKU\_PRASYARAT\_UNTUK \\
      LINTAS\_BUKU\_MEMPERDALAM \\
      LINTAS\_BUKU\_BERKAITAN\_DENGAN
    };

  % Klamp bounding box ke konten (target di kiri-atas, grade di
  % kanan-bawah) agar titik kontrol loop tidak menggelembungkan kotak.
  \pgfresetboundingbox
  \useasboundingbox
    ([xshift=-3mm,yshift=-3mm]xblabel.west |- target.south)
    rectangle
    ([xshift=3mm,yshift=3mm]d_bool.east |- grade.north);

\end{tikzpicture}%
  }
  \caption{Topologi Ontologi: Relasi Struktural, Keterurutan Penyajian
    (\texttt{NEXT\_CHAPTER}), dan Semantik (Lapisan Ekstraksi) Beserta
    Inferensi Lintas-Buku (Lapisan Completion).}
  \label{fig:ontologi}
\end{figure}

\subsubsection*{Definisi Kelas}

Ontologi terdiri dari lima kelas sebagaimana dirangkum pada
Tabel~\ref{tab:kelas-ontologi}.

\begin{table}[htbp]
  \centering
  \renewcommand{\arraystretch}{1.3}
  \caption{Definisi Kelas Ontologi}
  \label{tab:kelas-ontologi}
  \begin{tabularx}{\textwidth}{@{}lX@{}}
    \toprule
    \textbf{Kelas} & \textbf{Deskripsi} \\
    \midrule
    \texttt{Grade/Subject}
      & Tingkat kelas dan mata pelajaran pada Kurikulum Merdeka
        (mis.\ Kelas~XII Fisika). Akar hierarki. \\
    \texttt{Chapter}
      & Bab dalam buku teks, bersumber dari Daftar Isi.
        Memiliki ringkasan akademik (\texttt{summary}). \\
    \texttt{Subtopic}
      & Subbab di bawah \texttt{Chapter}.
        Mengelompokkan \texttt{Concept} dalam satu unit tematik. \\
    \texttt{Concept}
      & Konsep ilmiah yang diekstrak dari teks.
        Bersifat \textit{shared} lintas-dokumen melalui
        \texttt{MERGE} berdasarkan nama. \\
    \texttt{ConceptTarget}
      & Konsep yang disebutkan sebagai target relasi dalam ekstraksi
        per-bab tetapi tidak terdapat dalam himpunan konsep
        terekstrak resmi pada bab tersebut. Direpresentasikan sebagai
        kelas terpisah agar struktur relasi tetap terjaga tanpa
        mencampurkan \textit{stub} (yakni node konsep rintisan tanpa
        deskripsi) ke dalam populasi \texttt{Concept} formal. Tidak
        di-\textit{embed} (tidak dihitung representasi vektornya) dan tidak
        dijaring oleh tahap \textit{completion} lintas-buku. \\
    \bottomrule
  \end{tabularx}
\end{table}

\subsubsection*{Relasi Objek (Object Properties)}

Hierarki struktural mengikuti rantai:

\begin{center}
  \texttt{Grade/Subject} $\rightarrow$ \texttt{Chapter}
  $\rightarrow$ \texttt{Subtopic} $\rightarrow$ \texttt{Concept}
\end{center}

\noindent dihubungkan oleh relasi \texttt{HAS\_CHAPTER},
\texttt{HAS\_SUBTOPIC}, dan \texttt{HAS\_CONCEPT}. Pada tingkat
\texttt{Chapter} terdapat tiga relasi tambahan yang merekam
organisasi materi sebagaimana disusun dalam buku teks:
\texttt{NEXT\_CHAPTER} (urutan penyajian antar-bab),
\texttt{PRASYARAT} (ketergantungan konseptual antar-bab yang
terdokumentasikan dalam teks sumber, baik melalui rujukan langsung
maupun urutan penyajian materi), dan \texttt{MEMPERSIAPKAN} (relasi
inversi dari \texttt{PRASYARAT}). Ketiga relasi ini mengkodekan
keputusan editorial penyusun buku (yakni bagaimana materi disajikan
dalam sumber), bukan urutan pembelajaran yang dipreskripsikan
bagi siswa.

Vokabulari relasi semantik antarkonsep dibagi menjadi dua kelompok
dengan peran berbeda dalam \textit{pipeline}:

\begin{enumerate}
  \item \textbf{Relasi dalam-buku (\textit{within-book}):} lima belas tipe
        yang ditegakkan oleh \textit{prompt} LLM saat ekstraksi per-bab,
        menangkap struktur konseptual lokal dalam satu mata pelajaran
        (Tabel~\ref{tab:relasi-dalam-buku}).
  \item \textbf{Relasi lintas-buku (\texttt{LINTAS\_BUKU\_*}):} lima
        tipe yang dihasilkan pada tahap \textit{completion}
        lintas-disiplin, menangkap interkoneksi semantik antarmata
        pelajaran yang tidak tersurat dalam satu buku
        (Tabel~\ref{tab:relasi-lintas-buku}).
\end{enumerate}

\begin{table}[htbp]
\centering
\caption{Vokabulari Relasi Dalam-Buku
    (\texttt{Concept} $\leftrightarrow$ \texttt{Concept})}
\label{tab:relasi-dalam-buku}
\begin{tabular}{@{}p{0.32\linewidth}p{0.45\linewidth}p{0.16\linewidth}@{}}
  \toprule
  \textbf{Tipe Relasi} & \textbf{Semantik} & \textbf{Sifat} \\
  \midrule
  \texttt{BAGIAN\_DARI}
    & A merupakan bagian dari B (komposisi/hierarki)
    & asimetris, hierarkis \\
  \texttt{MENYEBABKAN}
    & A menyebabkan B (kausasi)
    & asimetris \\
  \texttt{BERGANTUNG\_PADA}
    & A bergantung pada B (dependensi fungsional)
    & asimetris \\
  \texttt{MENDEFINISIKAN}
    & A mendefinisikan B (definisional)
    & asimetris \\
  \texttt{MEMPENGARUHI}
    & A mempengaruhi B (pengaruh umum)
    & asimetris \\
  \texttt{BERINTERAKSI\_DENGAN}
    & A dan B berinteraksi resiprokal
    & simetris \\
  \texttt{MEMUNGKINKAN}
    & A memungkinkan B terjadi/berfungsi (\textit{enablement})
    & asimetris \\
  \texttt{MENGATUR}
    & A mengatur atau mengontrol perilaku B (regulasi)
    & asimetris \\
  \texttt{TERDIRI\_DARI}
    & A terdiri dari/tersusun atas B (komposisi/agregasi)
    & asimetris, hierarkis \\
  \texttt{BEREAKSI\_DENGAN}
    & A bereaksi dengan B (reaksi kimia)
    & simetris \\
  \texttt{MENGHASILKAN}
    & A menghasilkan B (produk atau keluaran proses)
    & asimetris \\
  \texttt{TERLETAK\_DI}
    & A terletak di/berada pada B (lokasi atau struktur)
    & asimetris \\
  \texttt{DIFORMULASIKAN\_SEBAGAI}
    & A diformulasikan sebagai B (representasi matematis/rumus)
    & asimetris \\
  \texttt{DIKATALISIS\_OLEH}
    & A dikatalisis oleh B (katalisis enzimatik/kimiawi)
    & asimetris \\
  \texttt{MEMILIKI\_SIFAT}
    & A memiliki sifat atau atribut B
    & asimetris \\
\end{tabular}
\end{table}

\begin{table}[htbp]
\centering
\caption{Vokabulari Relasi Lintas-Buku (Lintas-Disiplin)}
\label{tab:relasi-lintas-buku}
\begin{tabular}{@{}p{0.38\linewidth}p{0.42\linewidth}p{0.13\linewidth}@{}}
  \toprule
  \textbf{Tipe Relasi} & \textbf{Semantik} & \textbf{Sifat} \\
  \midrule
  \texttt{LINTAS\_BUKU\_SAMA\_DENGAN}
    & Konsep di buku berbeda merujuk pada entitas yang sama
    & simetris \\
  \texttt{LINTAS\_BUKU\_APLIKASI\_DARI}
    & Konsep A adalah penerapan konsep B di disiplin lain
    & asimetris \\
  \texttt{LINTAS\_BUKU\_PRASYARAT\_UNTUK}
    & Konsep A merupakan prasyarat pemahaman konsep B di mapel lain
    & asimetris, hierarkis \\
  \texttt{LINTAS\_BUKU\_MEMPERDALAM}
    & Konsep A memperdalam atau memperluas pemahaman konsep B
    & asimetris \\
  \texttt{LINTAS\_BUKU\_BERKAITAN\_DENGAN}
    & Berkaitan secara umum (\textit{fallback})
    & simetris \\
\end{tabular}
\end{table}

\subsubsection*{Properti Data (Data Properties)}

Setiap node dan edge dalam KG menyimpan properti data yang digunakan
untuk ekstraksi, validasi terminologi, dan komputasi
\textit{embedding}. Tabel~\ref{tab:data-properties} merangkum seluruh
properti beserta domain dan fungsinya.

\begin{table}[htbp]
  \centering
  \renewcommand{\arraystretch}{1.3}
  \caption{Properti Data Ontologi}
  \label{tab:data-properties}
  \begin{tabularx}{\textwidth}{@{}lll X@{}}
    \toprule
    \textbf{Properti} & \textbf{Tipe} & \textbf{Domain} & \textbf{Fungsi} \\
    \midrule
    \texttt{name}
      & \texttt{string} & Semua kelas
      & Nama tampilan node; kunci \texttt{MERGE} pada \texttt{Concept}. \\
    \texttt{grade}
      & \texttt{string} & Semua kelas
      & Atribut tingkat kelas, disalin untuk \textit{filtering} cepat. \\
    \texttt{type}
      & \texttt{string} & \texttt{Grade/Subject}
      & Tipe jenjang (mis.\ SMA, SMP). \\
    \texttt{summary}
      & \texttt{string} & \texttt{Chapter}
      & Ringkasan akademik 2--3 kalimat per bab. \\
    \texttt{chapter}
      & \texttt{string} & \texttt{Subtopic}
      & Nama \texttt{Chapter} induk (denormalisasi untuk kueri cepat). \\
    \texttt{description}
      & \texttt{string} & \texttt{Concept}
      & Deskripsi tekstual konsep; sumber utama \textit{embedding}
        untuk tahap KGC. \\
    \texttt{glossary\_validated}
      & \texttt{boolean} & \texttt{Concept}
      & Status validasi nama konsep terhadap glosarium buku teks. \\
    \texttt{materi\_pokok\_ref}
      & \texttt{string} & \texttt{Concept}
      & Referensi ke \textit{materi pokok} pada Buku Panduan Guru
        Kemendikbudristek sebagai jangkar penyelarasan kurikulum. \\
    \midrule
    \texttt{relation\_desc}
      & \texttt{string} & \textit{Edge}
      & Rasionalisasi satu kalimat (Bahasa Indonesia) atas eksistensi
        relasi; dihasilkan LLM dan disimpan bersama setiap edge. \\
    \texttt{relation\_type}
      & \texttt{string} & \textit{Edge}
      & Sub-tipe asli dari prediksi LLM untuk keperluan audit dan
        reprodusibilitas. \\
    \bottomrule
  \end{tabularx}
\end{table}
\end{document}
